\documentclass[10pt,a4paper]{amsart}
\usepackage{amsmath,amssymb,amsthm,graphicx}
\usepackage{verbatim}
\newenvironment{reference}{\comment}{\endcomment}
%\newenvironment{reference}{}{}
\newenvironment{slogan}{\comment}{\endcomment}
\newenvironment{history}{\comment}{\endcomment}

% For commutative diagrams we use Xy-pic
\usepackage[all]{xy}

% We use 2cell for 2-commutative diagrams.
\xyoption{2cell}
% Theorem environments.
%
\theoremstyle{plain}
\newtheorem{theorem}[subsection]{Theorem}
\newtheorem{proposition}[subsection]{Proposition}
\newtheorem{lemma}[subsection]{Lemma}

\theoremstyle{definition}
\newtheorem{definition}[subsection]{Definition}
\newtheorem{example}[subsection]{Example}
\newtheorem{exercise}[subsection]{Exercise}
\newtheorem{situation}[subsection]{Situation}

\theoremstyle{remark}
\newtheorem{remark}[subsection]{Remark}
\newtheorem{remarks}[subsection]{Remarks}

\numberwithin{equation}{subsection}

% Macros
%
\def\lim{\mathop{\mathrm{lim}}\nolimits}
\def\colim{\mathop{\mathrm{colim}}\nolimits}
\def\Spec{\mathop{\mathrm{Spec}}}
\def\Hom{\mathop{\mathrm{Hom}}\nolimits}
\def\Ext{\mathop{\mathrm{Ext}}\nolimits}
\def\SheafHom{\mathop{\mathcal{H}\!\mathit{om}}\nolimits}
\def\SheafExt{\mathop{\mathcal{E}\!\mathit{xt}}\nolimits}
\def\Sch{\mathit{Sch}}
\def\Mor{\mathop{\mathrm{Mor}}\nolimits}
\def\Ob{\mathop{\mathrm{Ob}}\nolimits}
\def\Sh{\mathop{\mathit{Sh}}\nolimits}
\def\NL{\mathop{N\!L}\nolimits}
\def\CH{\mathop{\mathrm{CH}}\nolimits}
\def\proetale{{pro\text{-}\acute{e}tale}}
\def\etale{{\acute{e}tale}}
\def\QCoh{\mathit{QCoh}}
\def\Ker{\mathop{\mathrm{Ker}}}
\def\Im{\mathop{\mathrm{Im}}}
\def\Coker{\mathop{\mathrm{Coker}}}
\def\Coim{\mathop{\mathrm{Coim}}}

% Boxtimes
%
\DeclareMathSymbol{\boxtimes}{\mathbin}{AMSa}{"02}

%
% Macros for moduli stacks/spaces
%
\def\QCohstack{\mathcal{QC}\!\mathit{oh}}
\def\Cohstack{\mathcal{C}\!\mathit{oh}}
\def\Spacesstack{\mathcal{S}\!\mathit{paces}}
\def\Quotfunctor{\mathrm{Quot}}
\def\Hilbfunctor{\mathrm{Hilb}}
\def\Curvesstack{\mathcal{C}\!\mathit{urves}}
\def\Polarizedstack{\mathcal{P}\!\mathit{olarized}}
\def\Complexesstack{\mathcal{C}\!\mathit{omplexes}}
% \Pic is the operator that assigns to X its picard group, usage \Pic(X)
% \Picardstack_{X/B} denotes the Picard stack of X over B
% \Picardfunctor_{X/B} denotes the Picard functor of X over B
\def\Pic{\mathop{\mathrm{Pic}}\nolimits}
\def\Picardstack{\mathcal{P}\!\mathit{ic}}
\def\Picardfunctor{\mathrm{Pic}}
\def\Deformationcategory{\mathcal{D}\!\mathit{ef}}
\title{A functorial acyclic replacement killing a prescribed set of complexes}
\author{The Stacks Project Authors}
\date{stacks-project:injectives.tex:local-snapshot-20261012; GNU FDL 1.2 or later}
\begin{document}
\maketitle
\noindent\textit{Editorial note.} All delivered mathematical text is the author's own, copied from the pinned Stacks Project chapter file sites-cohomology.tex; no statement or proof step was rewritten, reordered or removed.\par\medskip
\noindent\textit{Editorial note.} The proof is the author's complete construction, each of the two assertions being verified in the span; the closing sentence records that the second assertion holds by construction.\par\medskip
\noindent\textit{Editorial note.} Every labelled reference occurring inside the counted spans is carried in assembly.external\_references as a non-load-bearing navigation pointer (same-chapter ones at their definition site, cross-chapter ones at the invocation inside the counted span), so no load-bearing mathematics is replaced by a stub. No bibliography citation occurs inside the counted spans.\par\medskip
\noindent\textit{Editorial note.} The same-name check for this card scanned every staged Stacks chapter of this repository -- all sources/stacks-*/pinned-original-source.tex files, the card's own chapter excluded -- and the labels found are recorded in dedup.same\_name\_cross\_chapter on the card.\par\medskip
\noindent\textit{Editorial note.} The preamble is the chapter's own retained minimal preamble (sources/stacks-sites-cohomology/retained-preamble.tex), a verbatim excerpt of the native preamble: the Xy-pic setup and the native theorem declarations with the mathematical macros.\par\medskip
\noindent\textit{Allowed standard prerequisites (editorial).}\begin{itemize}
\item Functorial injective embeddings. Same chapter, Theorem theorem-injective-embedding-grothendieck. Supplies the functorial embeddings of the terms into injective objects out of which the replacement is built.
\item Products of injective objects and their mapping properties. Homology, the product rule for injective objects and the description of the relevant two-term complexes. Supply that products of the chosen injective objects are injective and that the maps out of the prescribed acyclic complexes into them vanish on homotopy, which is what makes the prescribed complexes redundant.
\end{itemize}\medskip
\begin{lemma}
\label{lemma-functorial-homotopies}
Let $\mathcal{A}$ be a Grothendieck abelian category.
Let $(K_i^\bullet)_{i \in I}$ be a set of acyclic complexes.
There exists a functor $M^\bullet \mapsto \mathbf{M}^\bullet(M^\bullet)$
and a natural transformation
$j_{M^\bullet} : M^\bullet \to \mathbf{M}^\bullet(M^\bullet)$
such
\begin{enumerate}
\item $j_{M^\bullet}$ is a (termwise) injective quasi-isomorphism, and
\item for every $i \in I$ and $w : K_i^\bullet \to M^\bullet$
the morphism $j_{M^\bullet} \circ w$ is homotopic to zero.
\end{enumerate}
\end{lemma}

\begin{proof}
For every $i \in I$ choose a (termwise) injective map of complexes
$K_i^\bullet \to L_i^\bullet$ which is homotopic to zero with
$L_i^\bullet$ quasi-isomorphic to zero. For example, take $L_i^\bullet$
to be the cone on the identity of $K_i^\bullet$.
We define $\mathbf{M}^\bullet(M^\bullet)$ by the following pushout diagram
$$
\xymatrix{
\bigoplus_{i \in I}
\bigoplus_{w : K_i^\bullet \to M^\bullet}
K_i^\bullet \ar[r] \ar[d] & M^\bullet \ar[d] \\
\bigoplus_{i \in I}
\bigoplus_{w : K_i^\bullet \to M^\bullet}
L_i^\bullet \ar[r] &  \mathbf{M}^\bullet(M^\bullet).
}
$$
Then $M^\bullet \mapsto \mathbf{M}^\bullet(M^\bullet)$ is a functor. The right
vertical arrow defines the functorial injective map $j_{M^\bullet}$.
The cokernel of $j_{M^\bullet}$ is isomorphic to the direct sum of
the cokernels of the maps $K_i^\bullet \to L_i^\bullet$ hence acyclic.
Thus $j_{M^\bullet}$ is a quasi-isomorphism. Part (2) holds by construction.
\end{proof}

\end{document}
